% ===========================================================================
%  user/06-live/00-chapter.tex
% ===========================================================================
\chapter{Live}\label{cap:live}
% A chapter full of boxes that do not break: with the page bottom aligned,
% all the spare space ended up between two paragraphs. Here the bottom stays
% ragged; at the end of the chapter the book's alignment comes back.
\raggedbottom

This is the page where a strategy stops being an exercise. A runner always
starts in simulation: it works out the orders on the connected account and
writes them to its journal, without sending them. Real orders are switched
on runner by runner, when you decide, with an explicit confirmation. From
then on, money moves.

\figuraSchermata{1.0}{figure/schermate/live.png}
  {The Live page. At the top the status, below it the three steps to create a
   runner, on the right the summary of what is missing and the saved
   runners.}
  {live}

\section{What a runner is}\label{sez:runner}

A \voce{runner} is what runs a strategy on an account. It is made of three
parts, and they are exactly what the three steps ask for:

\begin{enumerate}
  \item \textbf{A name.} It only serves to recognise it in the list.
  \item \textbf{What it runs.} A project (with its weights and its overlays),
        or a single strategy.
  \item \textbf{Where it runs.} Which broker. For MetaTrader 5 and
        Interactive Brokers there is also the choice of orders: simulated or
        real (\vediSez{real-orders}).
\end{enumerate}

\noindent
The panel on the right shows which of the three are in place: the
\pulsante{Save runner} button stays off until all three are. Below it, a line
says what the orders will be like, already in the colour they will have on
the runner's card.

\section{The three brokers}\label{sez:brokers}

\begin{center}
\begin{tabular}{@{}L{34mm}L{104mm}@{}}
\toprule
\textbf{Broker} & \textbf{What it does} \\
\midrule
\rigaTenue
\voce{TEST} & A simulated world with synthetic data and accelerated time.
  It needs no account and no credentials. It shows whether the machinery
  works --- retraining, composition, orders --- without waiting months. It
  is always simulated. \\
\voce{MetaTrader 5} & Runs on the account the MT5 terminal installed on the
  computer is logged into: Lyra needs no credentials. At start-up it checks
  that MT5 is there; with real orders, also that the account is
  \emph{netting}, that the broker allows automated trading and that
  \voce{Algo Trading} is on. \\
\rigaTenue
\voce{Interactive Brokers} & Connects to TWS or IB Gateway with the host,
  port, client ID and account saved with \pulsantesec{Connect}. If they are
  missing, it asks for them at start-up. \\
\bottomrule
\end{tabular}
\end{center}

\begin{pratica}{The order to follow, without exceptions}
\begin{enumerate}
  \item \textbf{TEST} for a few simulated days. It shows that the machinery
        starts, that orders come out, that the journals get written.
  \item \textbf{The real broker, in simulation.} The runner reads the
        account's equity and positions and writes to the journal the orders
        it would send. You find out which instruments the broker does not
        trade and which quantities do not add up.
  \item \textbf{A demo or paper account, with real orders on.} Real orders,
        play money: you find out about trading hours, rejections and
        missing permissions.
  \item \textbf{A real account}, with a fraction of the capital.
\end{enumerate}
Skipping a step saves no time: it only means finding the same problems later,
with the money in.
\end{pratica}

\section{Simulation or real orders}\label{sez:real-orders}

Once MetaTrader 5 or Interactive Brokers is chosen, the \voce{Orders} panel
appears below the broker cards, with two positions: \voce{Simulation}, the
starting one, and \voce{Real orders}.

\figuraSchermata{1.0}{figure/schermate/live-ordini.png}
  {A runner with real orders on: the \voce{Orders} panel and the summary line
   turn red. On the right, a runner already saved in simulation.}
  {live-orders}

\begin{pratica}{Switching real orders on}
\begin{enumerate}
  \item In step~\num{3}, choose \voce{MetaTrader 5} or
        \voce{Interactive Brokers}.
  \item In the \voce{Orders} panel, press \voce{Real orders}. A window opens
        that says what changes.
  \item Read it, tick the box and press \pulsante{Enable real orders}: the
        button stays off until the box is ticked.
  \item \pulsante{Save runner}. The \voce{REAL ORDERS} label appears in red on
        the runner's card.
\end{enumerate}
For a runner that is already saved, use \voce{Switch to real orders\ldots}
on its card, with the runner stopped: the window is the same.
\end{pratica}

\figuraSchermata{0.82}{figure/schermate/live-conferma.png}
  {The confirmation of real orders. The button turns on only after the box is
   ticked.}
  {live-confirm}

\begin{attenzione}{The runner manages the whole account}
It sizes the orders on the equity of the whole account, not on a part of
it, and in the symbols it trades it closes the positions that are not among
its targets: a share bought by hand on that account, if it is in the
strategy's universe, can be sold at the next cycle. That is why the account
must be dedicated to Lyra --- one per runner, no manual trading.
\end{attenzione}

\begin{nota}{First reduce, then increase}
In every cycle the runner first sends the orders that reduce positions,
then those that increase them; a reversal (from short to long, or the other
way round) becomes two orders, the close and the open. If part of the book
cannot go out --- old data on a stock, an order still open at the broker, a
holding with no price --- or a reduction does not go through, the orders
that increase exposure wait for the next cycle. A cycle left halfway never
leaves the account more exposed than before.
\end{nota}

\noindent
The mode is always visible: on the runner's card (\voce{SIMULATION} or
\voce{REAL ORDERS}) and on every line of the journal. It can only be changed
with the runner stopped, and going back to simulation with \voce{Back to
simulation} asks for no confirmation. The \pulsante{RUN} button of a runner
with real orders is red and asks once more before starting.

\begin{nota}{The confirmation stays on record}
The runner records when real orders were confirmed, and the journal reports
it. Going back to simulation deletes the confirmation: to switch them on
again, you confirm again.
\end{nota}

\subsection{The confirmation is for one account}\label{sez:bound-account}

At the first start after the confirmation, the runner records the account it
is connected to --- on MT5 the number and server, on IBKR the number --- and
from then on it does not move. If later the MT5 terminal is logged into
another account, or the \pulsantesec{Connect} settings point to another IBKR
account, the runner does not start and the journal says which account it
expected and which one it found. If it happens while it is running, it stops
before sending anything else. The account is shown on the runner's card and
in the start window.

\begin{pratica}{Moving to another account}
\begin{enumerate}
  \item With the runner stopped, \voce{Back to simulation}.
  \item Connect the broker to the new account (the MT5 terminal, or
        \pulsantesec{Connect} for IBKR).
  \item \voce{Switch to real orders\ldots} and confirm: at the next start the
        runner binds itself to the new account.
\end{enumerate}
\end{pratica}

\begin{nota}{Accounts in another currency}
Lyra sizes the orders in dollars, on the prices of its own data. An account
in euros is fine: the equity is converted with the exchange rate the broker
gives (on MT5, the EURUSD price in the terminal). If there is no rate, real
orders do not start and the journal says so.
\end{nota}

\section{The journal}\label{sez:journal}

Every runner keeps a journal, which opens with \voce{Journal} on its card.
Everything the runner does goes into it, one line per event, with the time
and the mode in which it happened:

\begin{itemize}
  \item starts and stops, with the account it connected to: number, type
        (demo, paper or real), terminal; with real orders, also the check
        that it is the confirmed account and Lyra's orders left open,
        cancelled at start and at stop;
  \item the orders. In simulation, those it would send; if they are the same
        at the next cycle, a single line says so. With real orders, those
        sent and their outcome: filled with quantity and price, partly
        filled, still working (\emph{pending}), cancelled, rejected with the
        reason the broker gives;
  \item warnings, for example a symbol the broker does not trade, or data
        that are too old;
  \item the kill-switch events (\vediSez{kill-switch});
  \item maintenance: retraining and data updates.
\end{itemize}

\figuraSchermata{1.0}{figure/schermate/live-registro.png}
  {The journal of a runner with real orders: two orders filled, one cancelled
   because it was not filled in time, a warning.}
  {live-journal}

\noindent
The filters at the top separate orders and fills, kill-switch, warnings and
errors. While the runner is running, the journal updates by itself.

\begin{nota}{It lives on the computer}
The journal is a file in the application folder,
\file{data/runners/<id>/journal.jsonl}, one line per event; beyond
\num{5}\,MB it continues in \file{journal.1.jsonl} and so on, up to ten
files. It stays after the application is closed and even after the runner is
deleted: it is the account's logbook, and it is worth keeping.
\end{nota}

\section{The kill-switch}\label{sez:kill-switch}

Three levels, on by default, the same as in the backtest:

\begin{center}
\begin{tabular}{@{}L{30mm}L{50mm}L{58mm}@{}}
\toprule
\textbf{Level} & \textbf{When it trips} & \textbf{What it does} \\
\midrule
\rigaTenue
Strategy & An \perc{8} drawdown & Suspends it for \num{30} days, then it
  comes back by itself. \\
Asset class & A \perc{5} loss in \num{30} days & Halves its exposure for
  \num{20} days. \\
\rigaTenue
Account & A \perc{20} drawdown from the peak & Closes the positions in the
  runner's symbols and opens no others. \\
\bottomrule
\end{tabular}
\end{center}

\noindent
A suspended strategy and a halved asset class stay that way even after the
runner is restarted, until they expire.

The account level, with real orders, belongs to the \textbf{account} and not
to the runner. The drawdown is measured at the start of every cycle, before
deciding: if the fall happened between one cycle and the next, the runner
does not send another round of purchases, but closes the positions in the
symbols it trades straight away. If something stays open (market closed, a
rejection) it tries again at once, then every minute three times, then at
the normal pace; once the positions are closed, it stops by itself.

\noindent
From then on the account stays closed: after a restart, after going to
simulation and back to real orders, for any runner with real orders on that
account. Starting a runner on a closed account only serves to close the
positions left: the start window says so, the runner closes and stops. While
it is closing, not even repeated errors stop it.

\begin{pratica}{Reopening the account}
\begin{enumerate}
  \item Stop the runner.
  \item \voce{Journal} \freccia \pulsantesec{Reset kill-switch}, and confirm.
        The account reopens for every runner that uses it.
  \item At the next start the peak starts again from the account's equity.
\end{enumerate}
\end{pratica}

\begin{perche}{Why it belongs to the account, and is saved}
The kill-switch protects money, and the money is in the account. When it
belonged to the runner, a restart was enough to start again from the equity
of the moment --- after a \perc{15} drawdown the \perc{20} limit in fact
became \perc{32} --- and going to simulation and back, or using another
runner, was enough to reopen a closed account. The state lives in
\file{data/accounts/}, one file per account; if it cannot be read, real
orders do not start: a closed account must not reopen because of a damaged
file.
\end{perche}

\begin{nota}{A fall found at start-up, withdrawals and deposits}
The kill-switch looks at the account's equity and does not tell a loss from
a withdrawal. If at start-up the equity is already below the account's peak
by more than the threshold --- the fall happened with the runner stopped ---
the kill-switch trips and is recorded, but the runner sends no orders and
says how to go on. If it is a loss, start the runner again: it closes the
positions left and stops. If money was withdrawn, reset the kill-switch from
the journal instead. If the account's equity cannot be read at start-up,
real orders do not start.
\end{nota}

\begin{nota}{In simulation}
The kill-switch trips in simulation too, and the journal says what it would
have done. There it closes nothing and does not touch the account: it
belongs to the runner, and it is reset the same way.
\end{nota}

\section{When orders do not go out}\label{sez:orders-held}

When an order does not go out, or does not go as expected, the journal says
so in one line. The most common ones, what they mean and what to do:

\begin{longtable}{@{}L{46mm}L{92mm}@{}}
\toprule
\textbf{In the journal} & \textbf{What it means, and what to do} \\
\midrule\endfirsthead
\toprule
\textbf{In the journal} & \textbf{What it means, and what to do} \\
\midrule\endhead
\rigaTenue
\enquote{orders NOT sent: the data are too old}
  & The prices the runner decides on are those downloaded in
    \percorso{Settings \freccia Data}, and they do not update by themselves.
    Every symbol must reach the last closed session of its market: for US
    shares and ETFs, the New York Stock Exchange calendar, holidays
    included; for forex and crypto, one session is tolerated. A day's bar is
    expected from the following day. The orders on the symbols that are
    behind wait, and so do those that increase exposure; reductions on the
    other symbols go out. Update the data: the held orders go out at the
    first cycle with fresh data. The closes ordered by the account's
    kill-switch go out regardless. \\
\enquote{real orders cannot start on this account}
  & The account cannot receive orders from Lyra, and the line says why: on
    MT5, a \emph{hedging} account (a \emph{netting} account is needed),
    trading disabled or not allowed to expert advisors, \voce{Algo Trading}
    off, the Python API switched off in the terminal's options; on IBKR, a
    login with several accounts and no account given in
    \pulsantesec{Connect}; on both, an account in another currency with no
    rate to convert it. \\
\rigaTenue
\enquote{the real orders of this runner were confirmed for \ldots{}}
  & The broker is connected to an account different from the one the runner
    started on after the confirmation. If it is the right account, go back
    to simulation and switch real orders on again (\vediSez{bound-account}). \\
\enquote{cycle skipped, nothing sent}
  & An obstacle that appeared while the runner was on, for example
    \voce{Algo Trading} switched off or the account data not yet in sync.
    The runner skips the cycles without sending anything and resumes by
    itself when it goes away. \\
\rigaTenue
\enquote{kill-switch: this account was closed}
  & The account was closed by the kill-switch: the runner starts only to
    close the positions left, then stops (\vediSez{kill-switch}). \\
\enquote{an order is still working at the broker}
  & On that symbol there is an order still open --- Lyra's or one placed by
    hand. While it stays open, the runner does not send another: when the
    first was filled, the position would double. \\
\rigaTenue
\enquote{order(s) that add exposure not sent}
  & Part of the book did not go out (one of the lines above, or a holding
    with no price in that cycle) or a sale did not go through: in that cycle
    only the reductions go out, and the purchases are worked out again at
    the next cycle. \\
An order \emph{rejected}
  & The reason is the broker's, written in the line: permissions, market
    closed, quantity below the minimum. In TWS and IB Gateway the
    \voce{Read-Only API} option must be off, otherwise every order is
    rejected. \\
\rigaTenue
An order \emph{cancelled} on IBKR
  & It was not filled within \num{30} seconds, usually because the market
    was closed, and the runner cancelled it: left open, it would have been
    filled later while the runner was working out another one. \\
An order \emph{pending}
  & The broker accepted it but has not filled it yet (MT5), or its
    cancellation is not yet confirmed (IBKR): it can still be filled. The
    runner sends no other orders on that symbol while it stays open, and
    cancels it when it stops. \\
\rigaTenue
An order \emph{skipped}
  & The quantity is below the broker's minimum lot, and nothing went out: the
    difference is smaller than one lot. \\
\enquote{the broker is not answering}
  & Connection dropped, TWS or MetaTrader~5 closed or restarting. The runner
    sends nothing and keeps trying without giving up, with waits that grow
    up to a quarter of an hour; when the broker answers again, it says so. \\
\rigaTenue
\enquote{cycle failed (1 of 3 in a row)}
  & An error of another kind. The runner tries again within a minute; three
    errors in a row stop it, with the cause in the journal --- except when it
    is closing an account closed by the kill-switch. \\
\enquote{cannot connect to TWS/IB Gateway}
  & TWS or IB Gateway off or not logged in, API not enabled
    (\percorso{Configure \freccia API \freccia Settings}), wrong port, or
    another program is using the same client ID. \\
\bottomrule
\end{longtable}

\section{The simulator}\label{sez:simulator}

The \voce{TEST SIMULATOR} panel at the top says how long a simulated day
lasts --- in the screenshot, two seconds --- and how much time it will cover
in all. Pressing \pulsantesec{Config} changes the pace.

\begin{nota}{What speeding up time is for}
A strategy that retrains every thirty days, on a real account, retrains for
the first time in thirty days. In the simulator it happens in a minute, and
if the retraining has a problem you find out at once instead of next month.
\end{nota}

\section{Credentials}\label{sez:credentials}

The broker credentials stay \textbf{on the computer of whoever uses the
application}, in a local file. They do not go through the site, they do not
end up in any database and they do not leave the machine.

For Interactive Brokers, \pulsantesec{Connect} saves the host, the port, the
client ID and the account. Port and client ID must be numbers, and the window
says so at once. The usual ports are \num{7497} for TWS in paper and
\num{7496} for real, \num{4002} and \num{4001} for IB Gateway. The account
must be given if the login manages more than one.

\begin{attenzione}{Two consequences}
The first: nobody can recover them for you: if they are lost, you set them
up again with the broker. The second, more important: whoever has access to
that computer has access to the account. The application is meant to run on
a personal machine, not on a server others can reach.
\end{attenzione}

\section{Starting and stopping}\label{sez:start-stop}

Once saved, the runner appears in the list at the bottom right, with the
commands to start and stop it. The two counters at the top ---
\voce{RUNNING} and \voce{STOPPED} --- say how many are in each state.

\figuraSchermata{1.0}{figure/schermate/live-runner.png}
  {Two runners on the same model: one stopped in simulation, one running
   with real orders.}
  {live-runner}

\noindent
The first cycle starts at once. The following ones follow the shortest bar
among those the runner watches, a quarter of a bar at a time: for daily
strategies, one cycle every six hours. Asking the broker more often would
produce identical cycles.

\noindent
With real orders, \pulsante{RUN} asks for confirmation once more and says
which account the orders will go to; if the account was closed by the
kill-switch, it says that starting will only serve to close positions.

\begin{nota}{Stopping: no new orders}
Once \pulsante{Stop} is pressed, the orders of the current cycle not yet sent
do not go out any more; those already sent reach their outcome, which goes
into the journal. With real orders, Lyra's orders still open at the broker
are cancelled: none of them is filled later without the runner watching.
\end{nota}

\begin{attenzione}{One broker per runner}
Two runners on the same account fight over the positions: the first opens,
the second sees a position it does not expect and closes it. The page says
so in step~\num{3}, and does not let a second runner start on the same
broker while the first is running.
\end{attenzione}

\flushbottom
